% id: 88-10
% book: 베이직쎈 확률과 통계 · 05 확률변수와 확률분포
% section: 유형 050 확률질량함수의 성질; 확률분포가 표로 주어진 경우
% figure: tikz:fig-88-10
% answer: 
% answer_source: 
% variant_level: 0 (원본 전사)
% 이 파일은 build.mjs 가 items.json 에서 만든다. 고칠 때는 items.json 을 고친다.
\smash{\raisebox{4.5mm}{\dmkichul{베이직쎈 확률과 통계 88쪽 10번}}}
확률변수 $X$의 확률분포를 표로 나타내면 다음과 같다. $\mathrm{P}(2\le X\le 3)=\dfrac{5}{12}$일 때, 상수 $a$, $b$에 대하여 $a-b$의 값은?
\par\smallskip\begin{center}\resizebox{\ifdim\width>\linewidth\linewidth\else\width\fi}{!}{% 88-10(인쇄 88쪽) — 확률변수 X의 확률분포표. 스캔 표라 tabular 로 다시 조판(칸 값 원문 그대로 · 답은 넣지 않는다).
{\setlength{\tabcolsep}{6pt}\renewcommand{\arraystretch}{1.8}%
\begin{tabular}{|c|c|c|c|c|c|}
\hline
$X$ & $0$ & $1$ & $2$ & $3$ & \textbf{합계} \\
\hline
$\mathrm{P}(X=x)$ & $a$ & $\dfrac{1}{6}$ & $b$ & $\dfrac{1}{3}$ & $1$ \\
\hline
\end{tabular}}
}\end{center}
\begin{choicesii}{\rule[-3ex]{0pt}{8ex}$\dfrac{1}{12}$}{\rule[-3ex]{0pt}{8ex}$\dfrac{1}{6}$}{\rule[-3ex]{0pt}{8ex}$\dfrac{1}{4}$}{\rule[-3ex]{0pt}{8ex}$\dfrac{1}{3}$}{\rule[-3ex]{0pt}{8ex}$\dfrac{5}{12}$}\end{choicesii}
